% =====================================================================
% Abbildung: ML-Node in der Sicherheitsarchitektur
% Passend zu Kapitel 18, Abschnitt "Qualität und Sicherheit"
%
% Einbindung im Buch, z. B. in chapters/18_ml_integration.tex:
%
%   \begin{figure}[H]
%     \centering
%     % =====================================================================
% Abbildung: ML-Node in der Sicherheitsarchitektur
% Passend zu Kapitel 18, Abschnitt "Qualität und Sicherheit"
%
% Einbindung im Buch, z. B. in chapters/18_ml_integration.tex:
%
%   \begin{figure}[H]
%     \centering
%     % =====================================================================
% Abbildung: ML-Node in der Sicherheitsarchitektur
% Passend zu Kapitel 18, Abschnitt "Qualität und Sicherheit"
%
% Einbindung im Buch, z. B. in chapters/18_ml_integration.tex:
%
%   \begin{figure}[H]
%     \centering
%     % =====================================================================
% Abbildung: ML-Node in der Sicherheitsarchitektur
% Passend zu Kapitel 18, Abschnitt "Qualität und Sicherheit"
%
% Einbindung im Buch, z. B. in chapters/18_ml_integration.tex:
%
%   \begin{figure}[H]
%     \centering
%     \input{figures/fig_kap18_ml_sicherheitsarchitektur}
%     \caption{Der ML-basierte \texttt{object\_detector\_node} liefert eine
%       unsichere Zusatzinformation, waehrend die Lidar-basierte
%       Hinderniserkennung die zuverlässigere Sicherheitsgrundlage bleibt.}
%     \label{fig:kap18-ml-sicherheitsarchitektur}
%   \end{figure}
%
% Benoetigte Pakete (im Buch bereits in main.tex geladen): tikz
% Benoetigte TikZ-Bibliotheken: positioning, arrows.meta, shadows, calc
% =====================================================================

\input{figures/fig_kap18_ml_sicherheitsarchitektur.tikz}

%     \caption{Der ML-basierte \texttt{object\_detector\_node} liefert eine
%       unsichere Zusatzinformation, waehrend die Lidar-basierte
%       Hinderniserkennung die zuverlässigere Sicherheitsgrundlage bleibt.}
%     \label{fig:kap18-ml-sicherheitsarchitektur}
%   \end{figure}
%
% Benoetigte Pakete (im Buch bereits in main.tex geladen): tikz
% Benoetigte TikZ-Bibliotheken: positioning, arrows.meta, shadows, calc
% =====================================================================

\begin{tikzpicture}[
    every node/.style={font=\small},
    node distance=10mm and 20mm,
    ros/.style={
        rectangle, rounded corners=1.5pt, draw=orange!60!black, thick,
        minimum width=42mm, minimum height=11mm,
        align=center, fill=orange!15, drop shadow
    },
    navnode/.style={ros, minimum width=36mm},
    flowunsafe/.style={-{Stealth[length=2.2mm]}, thick, black!60, dashed},
    flowsafe/.style={-{Stealth[length=2.2mm]}, thick, black!60},
    lbl/.style={font=\tiny, fill=white, inner sep=1pt, align=center}
]

% --- Pfad 1 (oben): ML ---
\node[ros] (mldet) at (0, 3) {\texttt{object\_detector\_node}\\(ML, Kamera)};

% --- Pfad 2 (unten): Lidar ---
\node[ros] (lidar) at (0, -3) {\texttt{lidar\_node} $\rightarrow$\\Hinderniserkennung};

% --- Gemeinsames Ziel ---
\node[navnode] (nav) at (8, 0) {\texttt{navigation\_node}};

\draw[flowunsafe] (mldet.east) -- node[lbl, above, text width=34mm]
    {Person erkannt?\\(unsicher)} (nav.north west);
\draw[flowsafe] (lidar.east) -- node[lbl, below, text width=34mm]
    {physische Distanzmessung\\(zuverlässiger)} (nav.south west);

% --- Hinweistext ---
\node[font=\small, align=center, text width=110mm, below=14mm of nav, xshift=-30mm]
    {ML ergänzt, ersetzt aber nicht die Lidar-basierte Sicherheitsprüfung.};

\end{tikzpicture}


%     \caption{Der ML-basierte \texttt{object\_detector\_node} liefert eine
%       unsichere Zusatzinformation, waehrend die Lidar-basierte
%       Hinderniserkennung die zuverlässigere Sicherheitsgrundlage bleibt.}
%     \label{fig:kap18-ml-sicherheitsarchitektur}
%   \end{figure}
%
% Benoetigte Pakete (im Buch bereits in main.tex geladen): tikz
% Benoetigte TikZ-Bibliotheken: positioning, arrows.meta, shadows, calc
% =====================================================================

\begin{tikzpicture}[
    every node/.style={font=\small},
    node distance=10mm and 20mm,
    ros/.style={
        rectangle, rounded corners=1.5pt, draw=orange!60!black, thick,
        minimum width=42mm, minimum height=11mm,
        align=center, fill=orange!15, drop shadow
    },
    navnode/.style={ros, minimum width=36mm},
    flowunsafe/.style={-{Stealth[length=2.2mm]}, thick, black!60, dashed},
    flowsafe/.style={-{Stealth[length=2.2mm]}, thick, black!60},
    lbl/.style={font=\tiny, fill=white, inner sep=1pt, align=center}
]

% --- Pfad 1 (oben): ML ---
\node[ros] (mldet) at (0, 3) {\texttt{object\_detector\_node}\\(ML, Kamera)};

% --- Pfad 2 (unten): Lidar ---
\node[ros] (lidar) at (0, -3) {\texttt{lidar\_node} $\rightarrow$\\Hinderniserkennung};

% --- Gemeinsames Ziel ---
\node[navnode] (nav) at (8, 0) {\texttt{navigation\_node}};

\draw[flowunsafe] (mldet.east) -- node[lbl, above, text width=34mm]
    {Person erkannt?\\(unsicher)} (nav.north west);
\draw[flowsafe] (lidar.east) -- node[lbl, below, text width=34mm]
    {physische Distanzmessung\\(zuverlässiger)} (nav.south west);

% --- Hinweistext ---
\node[font=\small, align=center, text width=110mm, below=14mm of nav, xshift=-30mm]
    {ML ergänzt, ersetzt aber nicht die Lidar-basierte Sicherheitsprüfung.};

\end{tikzpicture}


%     \caption{Der ML-basierte \texttt{object\_detector\_node} liefert eine
%       unsichere Zusatzinformation, waehrend die Lidar-basierte
%       Hinderniserkennung die zuverlässigere Sicherheitsgrundlage bleibt.}
%     \label{fig:kap18-ml-sicherheitsarchitektur}
%   \end{figure}
%
% Benoetigte Pakete (im Buch bereits in main.tex geladen): tikz
% Benoetigte TikZ-Bibliotheken: positioning, arrows.meta, shadows, calc
% =====================================================================

\begin{tikzpicture}[
    every node/.style={font=\small},
    node distance=10mm and 20mm,
    ros/.style={
        rectangle, rounded corners=1.5pt, draw=orange!60!black, thick,
        minimum width=42mm, minimum height=11mm,
        align=center, fill=orange!15, drop shadow
    },
    navnode/.style={ros, minimum width=36mm},
    flowunsafe/.style={-{Stealth[length=2.2mm]}, thick, black!60, dashed},
    flowsafe/.style={-{Stealth[length=2.2mm]}, thick, black!60},
    lbl/.style={font=\tiny, fill=white, inner sep=1pt, align=center}
]

% --- Pfad 1 (oben): ML ---
\node[ros] (mldet) at (0, 3) {\texttt{object\_detector\_node}\\(ML, Kamera)};

% --- Pfad 2 (unten): Lidar ---
\node[ros] (lidar) at (0, -3) {\texttt{lidar\_node} $\rightarrow$\\Hinderniserkennung};

% --- Gemeinsames Ziel ---
\node[navnode] (nav) at (8, 0) {\texttt{navigation\_node}};

\draw[flowunsafe] (mldet.east) -- node[lbl, above, text width=34mm]
    {Person erkannt?\\(unsicher)} (nav.north west);
\draw[flowsafe] (lidar.east) -- node[lbl, below, text width=34mm]
    {physische Distanzmessung\\(zuverlässiger)} (nav.south west);

% --- Hinweistext ---
\node[font=\small, align=center, text width=110mm, below=14mm of nav, xshift=-30mm]
    {ML ergänzt, ersetzt aber nicht die Lidar-basierte Sicherheitsprüfung.};

\end{tikzpicture}

